\section{Example Out-of-Loop Experiments in the nanoGPT Speedrun}
\label{app:labexp}
Each excerpt below is an experiment an agent created and ran outside the benchmark's training script during its nanoGPT run (\cref{sec:nanogpt}), reproduced from the traces and trimmed for length.

Kimi K3 re-derived Newton--Schulz iteration coefficients with a global optimizer, checked bf16 rounding bit-exactly:
\begin{lstlisting}[language=Python]
from scipy.optimize import differential_evolution
grid_in = np.concatenate([np.linspace(0.02, 0.05, 10),
                          np.linspace(0.05, 1.0, 190)])
grid_over = np.linspace(1.0, 1.3, 20)

def p_map(sig, a, b, c, iters=6):
    x = sig
    for _ in range(iters):
        x = a*x + b*x**3 + c*x**5
    return x

def objective(params):
    a, b, c = params
    dev = np.max(np.abs(p_map(grid_in, a, b, c) - 1.0))
    over = max(0.0, np.max(np.abs(p_map(grid_over, a, b, c))) - 1.15)
    return dev + 5.0*over

res = differential_evolution(objective,
        [(1.0, 6.0), (-8.0, 0.0), (0.0, 5.0)],
        maxiter=300, tol=1e-9, seed=0, polish=True)
\end{lstlisting}

DeepSeek V4 Pro built a calibrated toy of the training problem, with minibatch noise shaped by the true Kronecker Hessian and a natural-gradient oracle arm:
\begin{lstlisting}[language=Python]
"""Calibrated toy: Kron-quadratic + CORRECT Kron-Hessian
minibatch noise. eps = Hl^{1/2} Z Hr^{1/2} / sqrt(n_eff).
Ideal preconditioner: Ql = Hl^{-1/2}, Qr = Hr^{-1/2}."""
G = Hl @ W @ Hr
eps = Hl12 @ torch.randn(p, q) @ Hr12 / (n_eff ** 0.5)
g = G + eps
elif opt == "natgrad":   # oracle arm
    El, Vl = torch.linalg.eigh(Hl)
    Er, Vr = torch.linalg.eigh(Hr)
    u = Vl @ (Vl.T @ d @ Vr /
        (El[:, None]**0.5 * Er[None, :]**0.5)) @ Vr.T
\end{lstlisting}

GLM 5.3 debugged its SOAP implementation on CPU before any GPU screen:
\begin{lstlisting}[language=Python]
torch.manual_seed(0)
for shape in [(768, 768), (3072, 768), (768, 3072)]:
    p = torch.nn.Parameter((torch.randn(*shape) * 0.02).bfloat16())
    opt = SOAP([p], lr=0.025)
    for t in range(25):
        p.grad = (torch.randn(*shape) * 0.01).to(torch.bfloat16)
        opt.step()
        if not torch.isfinite(p.data).all():
            print(shape, 'NaN at step', t+1)
            st = opt.state[p]
            print('  L finite', torch.isfinite(st['L']).all().item(),
                  'v finite', torch.isfinite(st['v']).all().item())
            break
\end{lstlisting}

\section{Example Programmatic Orchestration}
\label{app:orchestration}
\begin{lstlisting}[language=Python]
# Admit independent subagents; do not wait for answers here.
review = await rlm("Audit the implementation. Reply with concrete issues.",
                   name="reviewer")
tests = await rlm("Run the test suite and classify failures.",
                  name="tester")

# Later, recover retained sessions and send a follow-up.
children = await rlm.list_subagents()
await agent_message.send(
    "Also inspect error-handling edge cases.",
    receiver_role="child", receiver_name=review.name)
\end{lstlisting}
The explicit reply path is intentional: a child is a persistent concurrent session, not a stateless completion returned by \rlmcall.

\section{LLM Usage Disclosure}
Large language models were used to assist with code development, writing refinement, and formatting during the preparation of this manuscript.
All scientific claims, experimental design, analysis, and intellectual contributions are solely the work of the authors.
